\documentclass[11pt,a4paper]{report}

\usepackage{float} %dzieki temu pakietowi rysunki pojawiaja sie tam gdzie chce po uzyciu opcji H w srodowisku figure
\usepackage[font=small,labelfont=bf]{caption}
\usepackage{natbib}
\usepackage{import}
\usepackage[tight]{subfigure}
\usepackage{tikz}
\usepackage{pgfplots}
\usepackage{polski}
\usepackage[cp1250]{inputenc}
\usepackage{indentfirst}
\usepackage{color}
\usepackage{graphics}
\usepackage{geometry}
% ---- PAGE LAYOUT ----
\geometry{a4paper, top=1.5cm, bottom=1.5cm, left=1.5cm,
right=1.5cm, nohead}
\usepackage{amsmath,amssymb,amsfonts}
\begin{document}
\begin{center}
\large{\textbf{PRACA DOMOWA II}}
\end{center}

\noindent
imi� i nazwisko ........................................\\


\noindent
\textbf{Zadanie 1} Wyznaczy� przedzia�y monotoniczno�ci podanych funkcji oraz ich ekstrema lokalne:\\
\textbf{(a)} $y = 3 x^4+x^3-2 x^2-x + 2$, 
\textbf{(b)} $y = (2x+3)^6(3x+2)^8$, 
\textbf{(c)} $y =\frac{x^2 - 3x + 4}{x - 3}$,  
\textbf{(d)} $y =  e^{-x^2} (2x+1)$.\\

\vspace{0.25cm}

\noindent\textbf{odpowiedzi}:\\ 
\textbf{(a)} $y\nearrow$ dla $x\in \left(-\frac{\sqrt{3}}{3},-\frac 14 \right) ,\ x\in \left(\frac{\sqrt{3}}{3},+\infty\right)$, 
$y\searrow$ dla $x\in \left(-\infty,-\frac{\sqrt{3}}{3}\right),\ x\in\left(-\frac 14,\frac{\sqrt{3}}{3}\right)$, maksimum lokalne dla $x = -\frac{\sqrt{3}}{3}$, minima lokalne dla $x = -\frac 14$, $x = \frac{\sqrt{3}}{3}$,\\ 
\textbf{(b)} $y\nearrow$ dla $x \in (-3/2,-8/7), x\in (-2/3,+\infty)$, 
$y\searrow$ dla $x\in (-\infty,-3/2), x\in (-7/8,-2/3)$, maksimum lokalne dla $x = -7/8$, minima lokalne dla $x = -3/2$, $x = -2/3$,\\ 
\textbf{(c)} $y\nearrow$ dla $x\in (-\infty,1),x\in(5,\infty)$, 
$y\searrow$ dla $x\in(1,3)$, $x\in(3,5)$, maksimum lokalne dla $x = 1$, minimum lokalne dla $x = 5$,\\ 
\textbf{(d)} $y\nearrow$ dla $x\in (-\infty, -1), x\in(1/2,\infty)$, $y\searrow$ dla $x\in(-1,1/2)$, maksimum lokalne dla $x = 1/2$, minimum lokalne dla $x = -1$.\\

\begin{figure}[H]
	\centering
	\begin{tikzpicture}
		\foreach \x in {0,...,34}
		\draw (0.5*\x,0)--(0.5*\x,18cm);
		\foreach \y in {0,...,36}
		\draw (0,0.5*\y)--(17cm,0.5*\y);
	\end{tikzpicture}
\end{figure}

\begin{figure}[H]
	\centering
	\begin{tikzpicture}
		\foreach \x in {0,...,34}
		\draw (0.5*\x,0)--(0.5*\x,27cm);
		\foreach \y in {0,...,54}
		\draw (0,0.5*\y)--(17cm,0.5*\y);
	\end{tikzpicture}
\end{figure}

\pagebreak


\noindent
\textbf{Zadanie 2} Oblicz ca�ki nieoznaczone:\\ 
\textbf{(a)} 1) $\int \left(\frac{5}{\cos^2 x} - 3 x^2 + \frac{1}{\sqrt[3]{x}}\right)dx$, 
2) $\int \frac{(2x-3)dx}{(x^2 - 3x + 5)^2}$, 3) $\int x \ln x dx$,\\
\textbf{(b)} 1) $\int \left(\frac{2}{x} + \frac{3}{x^2} - 4\sin x\right)dx$, 
2) $\int \frac{dx}{x (1-2\ln x)^3}$, 3) $\int x \cos(3x) dx$,\\
\textbf{(c)} 1) $\int \left(\frac{5}{\sqrt[4]{x^3}} - \frac{4}{x^2 + 1} + e^x\right)dx$, 
2) $\int \frac{\cos x dx}{\sqrt{1+3\sin x}}$, 3) $\int x\,\text{arctg}\,x dx$,\\
\textbf{(d)} 1) $\int \left(\frac{2}{\sqrt{1-x^2}} - 5\cos x + \frac{3}{\sqrt[3]{x^4}}\right)dx$, 
2) $\int \frac{\sin\left(\frac{1}{x^2}\right)dx}{x^3}$, 3) $\int \sin(3x)\sin(5x) dx$.\\


\vspace{0.25cm}
\noindent
\textbf{odpowiedzi}: \\
\textbf{(a)} 1) $5\,\text{tg}\,x - x^3 + \frac{3}{2}\sqrt[3]{x^2} + C$,
2) $-\frac{1}{x^2 -3x+5} + C$, 3) $\frac{x^2}{2}\ln x - \frac{x^2}{4} + C$,\\ 
\textbf{(b)} 1) $2\ln x - \frac{3}{x} + 4\cos x + C$,
2) $\frac{1}{4 (1 - 2 \ln x)^2} + C$, 3) $\frac{1}{3} x \sin (3 x)+\frac{1}{9} \cos (3 x) + C$,\\ 
\textbf{(c)} 1) $20\sqrt[4]{x}  - 4\,\text{arctg}\,x+e^x + C $,
2) $\frac{2}{3} \sqrt{3 \sin (x)+1} + C$, 3) $\frac{1}{2} x^2\,\text{arctg}\,x-\frac{x}{2}+\frac{1}{2} \,\text{arctg},x + C$,\\  
\textbf{(d)} 1) $2\arcsin x - 5\sin x - \frac{15}{\sqrt[3]{x}} + C$, 2) $\frac{1}{2} \cos\left(\frac{1}{x^2}\right) + C$, 3) $\frac{5}{16}\sin(3x)\sin(5x) + \frac{3}{16}\cos(3x)\cos(5x) + C$.\\

\begin{figure}[H]
	\centering
	\begin{tikzpicture}
		\foreach \x in {0,...,34}
		\draw (0.5*\x,0)--(0.5*\x,19cm);
		\foreach \y in {0,...,38}
		\draw (0,0.5*\y)--(17cm,0.5*\y);
	\end{tikzpicture}
\end{figure}

\begin{figure}[H]
	\centering
	\begin{tikzpicture}
		\foreach \x in {0,...,34}
		\draw (0.5*\x,0)--(0.5*\x,27cm);
		\foreach \y in {0,...,54}
		\draw (0,0.5*\y)--(17cm,0.5*\y);
	\end{tikzpicture}
\end{figure}

\pagebreak

\noindent
\textbf{Zadanie 3} Oblicz ca�ki nieoznaczone z funkcji wymiernych:\\ 
\textbf{(a)} $\int \frac{x+7}{x^2-x-6} dx$, 
\textbf{(b)} $\int \frac{-x^2-x-3}{x^3-3 x-2} dx$, 
\textbf{(c)} $\int \frac{4 x-16}{x^3-5 x^2+7 x+13} dx$, 
\textbf{(d)} $\int \frac{2 x^2+3 x-3}{x^3+3 x^2+7 x+5} dx$\\

\vspace{0.25cm}
\noindent
\textbf{odpowiedzi}: \\
\textbf{(a)} $2 \ln (3-x)-\ln (x+2) + C$,
\textbf{(b)} $-\frac{1}{x+1}-\ln (x-2) + C$, 
\textbf{(c)} $\frac{1}{2} \ln \left(x^2-6 x+13\right)-\ln (x+1) + C$,\\  
\textbf{(d)} $\frac{3}{2} \ln \left(x^2+2 x+5\right)-\ln (x+1)-\frac{1}{2} \,\text{arctg}\left(\frac{x+1}{2}\right) + C$.\\

\begin{figure}[H]
	\centering
	\begin{tikzpicture}
		\foreach \x in {0,...,34}
		\draw (0.5*\x,0)--(0.5*\x,22cm);
		\foreach \y in {0,...,44}
		\draw (0,0.5*\y)--(17cm,0.5*\y);
	\end{tikzpicture}
\end{figure}

\begin{figure}[H]
	\centering
	\begin{tikzpicture}
		\foreach \x in {0,...,34}
		\draw (0.5*\x,0)--(0.5*\x,27cm);
		\foreach \y in {0,...,54}
		\draw (0,0.5*\y)--(17cm,0.5*\y);
	\end{tikzpicture}
\end{figure}

\end{document}